\documentclass[11pt]{article}

\usepackage[margin=1in]{geometry}
\usepackage{booktabs}
\usepackage{graphicx}
\usepackage{hyperref}
\usepackage{microtype}
\usepackage{listings}
\lstset{basicstyle=\ttfamily\small, breaklines=true, frame=single}

\title{ThorAI: An Open-Source MCP-Enabled Benchmarking Platform for\\
NVIDIA DRIVE Thor}
\author{ThorAI Platform Contributors\\
\texttt{github.com/yourusername/thor-ai-platform}}
\date{Draft v0.1 --- \today}

\begin{document}
\maketitle

\begin{abstract}
NVIDIA DRIVE Thor is the automotive supercomputer that most developers
never get access to: it is expensive, scarce, and has no public
benchmark corpus. As a result, the community has no shared way to
compare model performance, power efficiency, or quantization gains
across the workloads that matter in a vehicle: real-time perception,
on-device LLM/VLM inference, and sensor fusion.

We present \textbf{ThorAI}, an open-source benchmarking and deployment
platform for NVIDIA DRIVE Thor. ThorAI contributes (1) the first
public benchmark suite for Thor with a unified result schema covering
latency, throughput, power, memory, and thermal metrics; (2) a Model
Context Protocol (MCP) server that lets AI assistants programmatically
run benchmarks, compare results, and manage models; (3) a TensorRT +
quantization optimization toolchain with executable INT8 dynamic
quantization; (4) reference implementations for bird's-eye-view (BEV)
sensor fusion and on-device vision-language models with automotive
safety filtering; and (5) community leaderboard infrastructure with a
moderated submission portal.

The platform is production-shaped: 5 Python packages, 107 passing
tests, PostgreSQL/InfluxDB storage with graceful in-memory fallback,
Docker deployment, and CI workflows. Reference measurements show
YOLOv8n at 71.3\,ms P50 latency with real GPU power/memory sampling,
and INT8 dynamic quantization achieving 2.0$\times$ weight compression,
all reproducible without a Thor device via deterministic simulation
mode.
\end{abstract}

\section{Introduction}
\label{sec:intro}

Autonomous driving stacks run a diverse set of models on
power-constrained embedded hardware: object detection, segmentation,
BEV perception, and increasingly large language and vision-language
models for scene understanding and driver assistance. NVIDIA DRIVE
Thor targets exactly this workload mix. Yet there is no open,
standardized way to answer basic questions that every deployment team
faces:
\begin{itemize}
  \item What latency/throughput/power does YOLOv8n achieve at batch 16?
  \item How much does INT8 quantization compress Llama-3-8B, and at
        what quality cost?
  \item Which precision profile meets a 50\,ms end-to-end budget for a
        given model?
\end{itemize}

ThorAI is our answer: an open platform that makes these measurements
first-class, machine-readable, and AI-assistable. The project has
three design pillars:
\begin{enumerate}
  \item \textbf{Schema-first results.} Every benchmark emits the same
        JSON result document (hardware, model, workload, metric
        sections), stored in PostgreSQL and streamed to InfluxDB.
  \item \textbf{MCP-native access.} Benchmarking is exposed as MCP
        tools (\texttt{benchmark.run}, \texttt{models.optimize},
        \texttt{benchmark.compare}, \ldots); any MCP-capable agent can
        drive the full benchmark--optimize--compare--report loop.
  \item \textbf{Honest, reproducible evaluation.} A deterministic
        simulation mode produces synthetic results for CI and
        development, while real runs collect live NVIDIA Management
        Library (NVML) telemetry: power, temperature, memory,
        utilization.
\end{enumerate}

\section{System Design}
\label{sec:design}

\subsection{Repository layout}
\begin{verbatim}
packages/
+-- core/
|   +-- thor-core/        # hardware monitor, metrics, config, timeseries
|   +-- thor-sdk/         # device abstraction, power, telemetry
|   +-- thor-mcp/         # MCP server: tools, resources, transports
+-- benchmarks/
    +-- thor-benchmark/   # runner, workloads, CLI
    +-- thor-models/      # registry, zoo, optimization toolchain
website/                  # leaderboard API (FastAPI) + React frontend
examples/                 # thor-sense (BEV), thor-vlm (VLM) references
tools/                    # Docker, CI, deployment scripts
\end{verbatim}

\subsection{Benchmark result schema}
Every run produces a single JSON document that doubles as the
leaderboard's database row:
\begin{center}
\begin{tabular}{ll}
\toprule
Section & Fields \\
\midrule
\texttt{hardware} & device, driver/CUDA/TensorRT versions, compute
capability, power limit, memory, utilization \\
\texttt{model} & name, source, architecture, parameters, precision,
input shape \\
\texttt{workload} & type, batch sizes, iterations, config \\
\texttt{results.latency} & P50/P95/P99/min/max/std (ms), count \\
\texttt{results.throughput} & samples/s, max batch size, tokens/s \\
\texttt{results.power} & average/peak/idle watts, joules/sample \\
\texttt{results.memory} & peak/average MB, allocation pattern \\
\texttt{results.thermal} & start/end/peak $^\circ$C, throttling events \\
\bottomrule
\end{tabular}
\end{center}

\subsection{Hardware monitoring}
\texttt{thor\_core.hardware.HardwareMonitor} samples GPU power,
temperature, memory, and utilization in a background thread during a
run using NVML (pynvml), with psutil for host telemetry. When no GPU
is available the monitor degrades gracefully and the affected sections
are marked \texttt{available: false}, keeping every run
schema-complete.

\subsection{Workloads}
\begin{center}
\begin{tabular}{ll}
\toprule
Workload type & Models (built-in zoo) \\
\midrule
\texttt{vision} (detection) & YOLOv8n/s/m/l, DETR-ResNet50, RT-DETR \\
\texttt{segmentation} & YOLOv8n-seg, SegFormer-B0 \\
\texttt{classification} & ResNet-50, ViT-Base \\
\texttt{language} & Llama-3-8B, Mistral-7B, Phi-3-mini, Qwen2-7B,
Gemma-7B \\
\texttt{multimodal} & LLaVA-1.5-7B, Qwen-VL-Chat \\
\bottomrule
\end{tabular}
\end{center}
Model loading is lazy (torch/ultralytics/transformers); a
\texttt{simulate} flag produces deterministic synthetic results for CI
and development.

\subsection{MCP server}
The MCP server exposes 13 tools, 4 resources, and 3 prompts:
\begin{itemize}
  \item \textbf{Tools}: \texttt{benchmark\_run},
        \texttt{benchmark\_compare}, \texttt{benchmark\_list},
        \texttt{models\_list}, \texttt{models\_register},
        \texttt{models\_optimize}, \texttt{models\_deploy},
        \texttt{datasets\_list}, \texttt{datasets\_register},
        \texttt{reports\_generate}, \texttt{hardware\_status},
        \texttt{experiments\_track}, \texttt{experiments\_list}.
  \item \textbf{Resources}: \texttt{thor://benchmarks/results},
        \texttt{thor://models/registry},
        \texttt{thor://hardware/telemetry},
        \texttt{thor://experiments/history}.
  \item \textbf{Transports}: stdio (default), a FastAPI REST bridge,
        and a standard streamable-HTTP MCP endpoint at
        \texttt{/mcp} so the server can be hosted as a remote MCP
        server; verified interoperable against a live third-party
        remote MCP endpoint.
\end{itemize}
Auth uses HMAC-signed bearer tokens; rate limiting is a per-key token
bucket. Storage uses PostgreSQL when configured (including via
\texttt{DATABASE\_URL} from deployment platforms) with automatic
in-memory fallback.

\subsection{Optimization toolchain}
\texttt{models\_optimize} creates optimization profiles (TensorRT,
quantization, pruning, distillation). Execution is real where the
toolchain is present:
\begin{itemize}
  \item \textbf{INT8 quantization} --- executable torch dynamic
        quantization of \texttt{nn.Linear}/\texttt{nn.LSTM} weights
        with size/compression reporting (verified: 2.0$\times$ weight
        compression on a reference model).
  \item \textbf{TensorRT} --- torch ONNX export to a TensorRT engine
        with min/opt/max batch optimization profiles, FP16/INT8
        flags, an INT8 calibration hook, and \texttt{.plan}
        serialization. Requires the TensorRT toolchain on-device; the
        build logic is unit-tested against a mock TensorRT module.
\end{itemize}
INT4 (GPTQ-style) and FP8 execution are staged pending the
bitsandbytes/GPTQ toolchain.

\subsection{Reference implementations}
\begin{itemize}
  \item \textbf{thor-sense (BEV)}: camera and LiDAR encoders, BEV
        projection and fusion modules, pinhole projection,
        BEV-IoU-based detection fusion, and a constant-velocity
        object tracker.
  \item \textbf{thor-vlm}: a vision encoder + projector + tiny causal
        language model reference stack running end-to-end on CPU, a
        transformers backend for real VLMs (e.g.\ LLaVA), and
        automotive safety filters (prompt/output denylists) applied to
        every generation.
\end{itemize}

\section{Evaluation}
\label{sec:eval}

\subsection{Software correctness}
The platform ships \textbf{107 passing tests} across 5 packages and 3
example suites, covering metric correctness, hardware fallback,
config parsing, MCP tool dispatch, the streamable-HTTP transport
(end-to-end against a live uvicorn server), quantization execution,
TensorRT build logic, sensor fusion, VLM generation, and the
submission portal API. The frontend builds cleanly under TypeScript
strict mode.

\subsection{Reference measurements (dev workstation)}
Measured on a development workstation (RTX 3050 Ti Laptop GPU, CPU
torch inference):
\begin{center}
\begin{tabular}{lcccc}
\toprule
Model & Workload & P50 latency & Throughput & Power (avg) \\
\midrule
YOLOv8n & detection & 71.3\,ms & 13.9 samples/s & 7.7\,W \\
INT8-quantized reference model & --- & --- & --- & 2.0$\times$ compression \\
\bottomrule
\end{tabular}
\end{center}
Real GPU power/memory sampling was exercised during runs
(7.7\,W average, $\sim$430\,MB peak memory), validating the telemetry
pipeline end to end. Thor-device measurements (TensorRT engines,
FP16/INT4 LLM decode, BEVFormer) are collected per the runbook and
will be reported in the camera-ready version.

\subsection{Reproducibility}
Every result is a schema-complete JSON document with hardware and
environment context, stored in PostgreSQL, streamed to InfluxDB, and
aggregatable into a community leaderboard. CI workflows re-run
benchmarks on push and enforce regression thresholds (latency $+10\%$,
throughput $-5\%$, power $+15\%$).

\section{Related Work}
\label{sec:related}
Generic inference benchmarks (e.g.\ MLPerf) and vendor SDK demos do
not target Thor, do not expose results through MCP, and do not ship
reference automotive model implementations. To our knowledge ThorAI
is the first open platform combining a public Thor benchmark suite,
MCP-driven benchmarking, and reference automotive models.

\section{Conclusion and Future Work}
\label{sec:conclusion}
We presented ThorAI, an open-source, MCP-enabled benchmarking platform
for NVIDIA DRIVE Thor, spanning benchmarking, optimization, and
community leaderboard infrastructure. The platform is tested (107
tests), deployable (Docker, remote MCP hosting, CI), and reproducible
(deterministic simulation + schema-complete results).

Future work: (1) a public Thor benchmark corpus with leaderboard
submissions; (2) deep BEVFormer-style inverse-projection and
pretrained VLM fine-tuning; (3) INT4/FP8 quantization execution; (4)
multi-GPU and distributed inference benchmarks.

\begin{thebibliography}{9}
\bibitem{mcp}
Model Context Protocol specification. \url{https://modelcontextprotocol.io}
\bibitem{thor}
NVIDIA DRIVE Thor. \url{https://www.nvidia.com/en-us/self-driving-cars/drive-thor/}
\bibitem{tensorrt}
NVIDIA TensorRT. \url{https://developer.nvidia.com/tensorrt}
\bibitem{mlperf}
MLPerf Inference benchmark. \url{https://mlcommons.org/benchmarks/inference/}
\end{thebibliography}

\end{document}
